\chapter{The four-tank protocol (version 2)}
\label{app:four-tank-protocol}
% Moved from the former coordination chapter (sec:coord-protocol). Registered
% in the thesis repository (results/coordination/PROTOCOL.md) and never run
% (Section sec:tank-protocol). Kept as registered; shorten if pages are tight.
During the exploratory coordination runs, the training loop was redesigned
about a dozen times while the development, held-out and beyond results were
being looked at, and the best champion came out of a hand-steered chain of
runs. Such a result cannot test a hypothesis. Before any confirmatory run,
a protocol was therefore registered in the repository. It fixes, in one
commit, the test bed, the score weights, the batteries, the method, the
number of runs, the comparator and the statistics. Any change creates a new
version and restarts the runs; runs under an earlier version are reported,
not discarded. Version~2, the current one, was registered after the PI loops
of the test bed were retuned and the MPC re-tuned accordingly, before any
protocol run and before any access to the sealed battery.

\paragraph{Runs and policies.} $N = 8$ independent runs of six generations
each, all from the fixed recipe, with \texttt{deepseek-flash} at high
reasoning effort and three candidates per generation. A run that fails for
technical reasons is resumed under the same name, never replaced or dropped.
The policy of each run is its final champion; it is not chosen by any other
score.

\paragraph{Comparator.} The primary comparator is the tuned MPC with frozen
parameters; the untuned MPC is secondary, and the fixed recipe and the tuned
MPC with known disturbance are references.

\paragraph{Sealed test battery.} 183 scenarios with their own seed: 42 in the
development range (including nominal target changes), 36 each at 1.25, 1.5
and 2 times a development draw, 24 with sine periods outside the development
range (100--180~s and 550--800~s), and 9 with feed, pump and split
disturbances together. The battery is evaluated once, after all runs are
complete.

\paragraph{Analyses.} All intervals come from a hierarchical bootstrap
($B = 10\,000$) that resamples policies and then scenarios.
\begin{itemize}
\item $H_{0,1}$: the mean paired difference in score, LLM minus tuned MPC,
  over the 8 policies and the 42 development-range test scenarios. $H_{0,1}$
  is rejected if the 95\% interval excludes zero. Failing to reject does not
  establish parity: the two are called equivalent only if the 90\% interval
  lies within $\pm 10\%$ of the MPC's mean score (two one-sided tests).
\item $H_{0,2}$: the slope of the mean $\log(1 + J_s)$ against the amplitude
  level (1, 1.25, 1.5 and 2), LLM minus MPC. $H_{0,2}$ is rejected if the
  95\% interval of the slope difference excludes zero.
\item C3: the median, interquartile range and range of the per-policy mean
  scores, the between-policy standard deviation with a chi-square
  confidence interval, and the number of runs without any promotion.
\end{itemize}
Secondary analyses are the per-term differences with Holm-corrected
bootstrap $p$-values, the score ratios on the unseen periods and
combinations, and results per disturbance type and pattern. Stress tests
with noise, plant mismatch, analyser delay and telemetry faults, and two
ablations (a prompt without the plant equations, and the best of 18 one-shot
samples without the feedback loop) are planned as exploratory.
